\documentclass[pdflatex,sn-mathphys-num]{sn-jnl}

\usepackage{graphicx}
\usepackage{amsmath,amssymb,amsfonts}
\usepackage{amsthm}
\usepackage{mathtools}
\usepackage{xcolor}
\usepackage{booktabs}
\usepackage{listings}

\lstset{
  basicstyle=\footnotesize\ttfamily,
  breaklines=true,
  breakatwhitespace=true,
  frame=single,
  framerule=0.4pt,
  rulecolor=\color{black!20}
}

\theoremstyle{thmstyleone}
\newtheorem{theorem}{Theorem}[section]
\newtheorem{lemma}[theorem]{Lemma}
\newtheorem{proposition}[theorem]{Proposition}
\newtheorem{corollary}[theorem]{Corollary}
\newtheorem{conjecture}[theorem]{Conjecture}
\theoremstyle{thmstyletwo}
\newtheorem{example}[theorem]{Example}
\newtheorem{remark}[theorem]{Remark}
\theoremstyle{thmstylethree}
\newtheorem{definition}[theorem]{Definition}

\DeclareMathOperator{\Aut}{Aut}
\DeclareMathOperator{\Fix}{Fix}
\DeclareMathOperator{\Tr}{Tr}
\DeclareMathOperator{\Spec}{Spec}
\DeclareMathOperator{\Stab}{Stab}
\DeclareMathOperator{\Sym}{Sym}

\numberwithin{equation}{section}

\raggedbottom

\begin{document}

\title[A Kernel-Checked Order-7 Exclusion for Conway's 99-Graph]{A Kernel-Checked Lean~4 Proof That Conway's 99-Graph Has No Automorphism of Order~7, with Certified SAT Cross-Checks}

\author*[1]{\fnm{Antonio} \sur{Machuca}}\email{am.machuca.2023@alumnos.urjc.es}

\affil*[1]{\orgdiv{Department of Computer Science and Statistics}, \orgname{Universidad Rey Juan Carlos}, \orgaddress{\city{Madrid}, \country{Spain}}}

\abstract{Conway's 99-graph problem asks if a strongly regular graph with parameters $(99,14,1,2)$ exists. The problem is open.
Behbahani and Lam used a computer search to show that such a graph has no automorphism of order~7.
This article gives a proof of that result that the Lean~4 kernel checks.
The theorem \texttt{conway\_no\_z7\_automorphism} has no hypotheses other than the graph conditions and the automorphism.
For this theorem, \texttt{\#print axioms} reports only \texttt{[propext, Quot.sound]}.
The proof uses no SAT solver, no external certificate, and no \texttt{native\_decide}.
The proof changes an automorphism of order~7 into a $12 \times 12$ system of integer equations.
A change of labels puts the first row of each solution in a list of 33 rows.
Short arithmetic arguments remove 5 of these rows, and finite searches in the kernel remove the other 28 rows.
Each kernel check stays below 3\,GB of memory.
For this purpose, the library stores row lists that are the same for all 28 rows as checked theorems. It also divides each large search into many small theorems.
Two CNF refutations, each checked by \texttt{drat-trim}, give independent checks of the order-7 case at the formula level.
The article also records kernel-checked structural lemmas, a repaired defect in a SAT encoder, and positive controls. A solver campaign on the two open symmetric cases gave no result.
Four labels (\textsc{Proved}, \textsc{Compiled}, \textsc{Explored}, \textsc{Pending}) give the status of each claim.
The article does not decide if the graph exists.}

\keywords{Conway 99-graph, strongly regular graphs, Lean 4, interactive theorem proving, kernel computation, DRAT certification}

\maketitle

\section{Introduction}

\subsection{The problem}

A strongly regular graph $\mathrm{srg}(v, k, \lambda, \mu)$ is a $k$-regular graph on $v$ vertices.
Two adjacent vertices have exactly $\lambda$ common neighbors.
Two distinct vertices that are not adjacent have exactly $\mu$ common neighbors.
It is not known if a strongly regular graph with parameters $(99, 14, 1, 2)$ exists~\cite{brouwer2011spectra,cesarz2025automorphisms}.
Conway put this question in a list of five problems with a prize of \$1{,}000.
The text in circulation is the prize list \emph{Five \$1,000 Problems (Update 2017)}~\cite{conway1970}, Problem~2.
Bill Cheswick added an update in 2014, and Cheswick and Sloane added a postscript in 2017.
That document is not a research article from 1970.

Let $A \in \{0,1\}^{99 \times 99}$ be the adjacency matrix of a Conway 99-graph. Then
\begin{equation}
 A = A^T, \qquad \operatorname{diag}(A) = 0, \qquad A^2 + A = 12I_{99} + 2J_{99},
 \label{eq:matrix}
\end{equation}
where $I_n$ is the identity matrix and $J_n$ is the all-ones matrix of size $n \times n$.
The spectrum of $A$ is
\begin{equation}
 \Spec(A) = \{14^1, 3^{54}, (-4)^{44}\}.
 \label{eq:spectrum}
\end{equation}
For a vertex $x$, let $\Gamma(x)$ be its set of neighbors. Let $\Gamma_2(x)$ be the set of the $84$ vertices that are not $x$ and not in $\Gamma(x)$.
Because $\lambda = 1$, each vertex of $\Gamma(x)$ has exactly one neighbor in $\Gamma(x)$.
As a result, $\Gamma(x)$ induces seven disjoint edges ($7K_2$).
Because $\mu = 2$, each vertex of $\Gamma_2(x)$ has exactly two neighbors in $\Gamma(x)$.

\subsection{Results from the literature}

A search of all graphs on 99 vertices has $2^{\binom{99}{2}} = 2^{4851}$ edge sets. For this reason, many searches add a nontrivial automorphism. Such a search cannot find a rigid graph. Let $G$ be a Conway 99-graph. The literature gives these results for $\Aut(G)$.
\begin{enumerate}
 \item \textbf{An involution fixes exactly one vertex.} Makhnev and Minakova~\cite{makhnev2001automorphisms} classify the fixed subgraph of an automorphism of prime order. Makhnev's 2009 slides~\cite{makhnev2009symmetric} state the result for involutions and give it to Makhnev and Minakova.
 \item \textbf{Only the primes 2 and 3 divide $|\Aut(G)|$, and an automorphism of order~3 has no fixed points.} Behbahani and Lam prove both statements~\cite{behbahani2009thesis,behbahani2011strongly}. In the thesis, the statement is Theorem~4.14. Their orbit-matrix search for order~3 without fixed points did not finish. Table~21 of the thesis shows ``?'' for that case~\cite{behbahani2009thesis}.
 \item \textbf{No group of order 6 or 9 acts.} Crnkovi{\'c} and Maksimovi{\'c}~\cite{crnkovic2014strongly} exclude automorphism groups of order 6 and order 9. They cite the fixed-point result for order~3 from Behbahani and Lam. They do not exclude an automorphism of order~3 without fixed points.
 \item \textbf{Computer-free proofs.} Cesarz and Woldar~\cite{cesarz2025automorphisms} prove these two statements without a computer:
 \begin{equation}
  2 \mid |\Aut(G)| \implies |\Aut(G)| \mid 6,
  \qquad
  7 \mid |\Aut(G)| \implies \Aut(G) \cong \mathbb{Z}_7.
  \label{eq:prior_bounds}
 \end{equation}
 \item \textbf{CP-SAT search.} Thakkar and Severini~\cite{thakkar2026forced} encode forced structure and a prescribed automorphism in Google OR-Tools CP-SAT. Their order-7 instance with one fixed vertex returned \texttt{UNKNOWN} after 48 hours on 14 cores.
\end{enumerate}
Items 2 to 4 together give $|\Aut(G)| \in \{1, 2, 3\}$.
If $|\Aut(G)|$ is even, then $|\Aut(G)|$ divides 6, and 6 is excluded.
If $|\Aut(G)|$ is odd, then $|\Aut(G)| = 3^b$, and a group of order $3^b$ with $b \ge 2$ has a subgroup of order 9.
As a result, two symmetric cases stay open: an involution that fixes exactly one vertex, and an automorphism of order~3 without fixed points.
This article does not formalize this consequence.

\subsection{Contributions}

\begin{itemize}
 \item \textbf{A kernel-checked proof of the order-7 exclusion (\textsc{Proved}).} Theorem~\ref{thm:no_z7} is \texttt{conway\_no\_z7\_automorphism}. It states that no matrix that satisfies the graph conditions has a nontrivial permutation $s$ with $s^7 = \mathrm{id}$ that keeps the adjacency. The axiom report is \texttt{[propext, Quot.sound]}. The result is the known theorem of Behbahani and Lam~\cite{behbahani2011strongly}. The contribution is a machine-checked proof that does not depend on a SAT solver or on an external certificate.
 \item \textbf{A method for large kernel computations.} The proof contains 28 finite searches. Section~\ref{sec:resources} describes how we keep each kernel check below 3\,GB of memory. It also gives a measurement error that did not show the real memory use.
 \item \textbf{Kernel-checked structural lemmas (\textsc{Proved}).} The library proves $K_4$-freeness, the $7K_2$ structure of each neighborhood, diameter at most 2, and a checker that is equivalent to the graph conditions. These facts are standard.
 \item \textbf{Two certified CNF refutations (\textsc{Proved} as formulas).} \texttt{drat-trim} verifies a refutation of a canonical order-7 CNF and of a CNF for the exterior equations. The Lean proof does not use them. The step from the graph to each CNF is \textsc{Compiled}.
 \item \textbf{Encoder validation and open cases.} We record a repaired defect in a SAT encoder, positive controls, and a column-weight lemma for the open involution case. A solver campaign on the two open cases gave no result (\textsc{Explored}).
\end{itemize}

\section{Status labels}
\label{sec:taxonomy}

Each result in this article has one of four labels.
\begin{description}
 \item[\textsc{Proved}.] The Lean~4 kernel checks the theorem, and \texttt{\#print axioms} reports only \texttt{[propext, Quot.sound]}. For a SAT claim, the label applies to one CNF file. The solver log of that file must end in \texttt{s UNSATISFIABLE}, and \texttt{drat-trim} must report \texttt{s VERIFIED} for that file and its proof.
 \item[\textsc{Compiled}.] The Lean code compiles, but a theorem uses \texttt{sorry} or a premise that the kernel does not derive. The label also applies to a SAT compiler that passes deterministic regression tests.
 \item[\textsc{Explored}.] A search or a calculation finished without a satisfying assignment and without a checked refutation.
 \item[\textsc{Pending}.] An open mathematical question, or a formal link that is not done.
\end{description}
A \textsc{Proved} label applies to the stated theorem or file only. It does not apply to an interpretation that the text adds.

\section{Lean definitions and structural lemmas}

\subsection{The graph conditions and a checker}

The library uses Lean~4.33.1~\cite{moura2021lean4} and no external Lean library.
The predicate \texttt{ConwayAdj} states that $A$ is symmetric, has a zero diagonal and binary entries, and satisfies~\eqref{eq:matrix}.
The Boolean function \texttt{checkConway} computes these conditions. Two theorems give the equivalence:
\begin{equation}
 \texttt{conway\_soundness}, \texttt{conway\_complete}: \quad \forall A \in \mathbb{N}^{99 \times 99}, \quad \texttt{checkConway}(A) = \texttt{true} \iff \texttt{ConwayAdj}(A).
 \label{eq:checker}
\end{equation}

\subsection{Structural lemmas (\textsc{Proved})}

The module \texttt{Structural.lean} proves these lemmas from \texttt{ConwayAdj}:
\begin{itemize}
 \item \texttt{conway\_k4\_free}: the graph contains no $K_4$.
 \item \texttt{conway\_neighborhood\_one\_factor}: each neighborhood induces $7K_2$.
 \item \texttt{conway\_diameter\_at\_most\_two}: two vertices that are not adjacent have a common neighbor.
\end{itemize}
The module \texttt{Z2Classification.lean} proves \texttt{conway\_z2\_f3\_fixed\_points\_dichotomy}. If an involution fixes three vertices, then these vertices induce $K_3$ or $3K_1$.
For all 58 declarations in \texttt{Structural.lean}, \texttt{\#print axioms} reports only \texttt{[propext, Quot.sound]}.
Two earlier proofs used \texttt{Classical.choice}. The source text did not show it. The axiom report showed it.
The causes were the library lemma \texttt{beq\_self\_eq\_true} and a call of \texttt{omega} on a branch with no cases. We replaced both steps by explicit case analysis. The statements did not change.

\subsection{Arithmetic refutations and graph premises}

Character-theoretic arguments give numerical conditions on fixed edges and eigenspace multiplicities. The library proves the arithmetic refutations of these conditions:
\begin{itemize}
 \item For an involution with $f=5$ fixed vertices, \texttt{conway\_z2\_f5\_\allowbreak spectral\_\allowbreak formula\_\allowbreak contradiction} shows that no integer $\varepsilon_1$ satisfies
 \begin{equation}
  (\varepsilon_1 = 18 \lor \varepsilon_1 = 15) \land (\varepsilon_1 \bmod 7 = 6).
  \label{eq:f5_arith}
 \end{equation}
 \item For $f=7$, \texttt{conway\_no\_z2\_f7\_\allowbreak formula\_\allowbreak automorphism} shows that no integer satisfies
 \begin{equation}
  \varepsilon_1 \in \{7, 10, 13\} \land (\varepsilon_1 \bmod 7 = 2).
  \label{eq:f7_arith}
 \end{equation}
 \item \texttt{CesarzWoldarTheorems.lean} proves final arithmetic contradictions, for example $7a = 62$ or a sum of even numbers that is equal to 5.
\end{itemize}
These arithmetic refutations are \textsc{Proved}.
Their graph-level wrappers, for example \texttt{conway\_no\_z2\_f5\_\allowbreak automorphism}, take the counting identities and the spectral congruences as hypotheses. These wrappers are \textsc{Compiled}.
A derivation of these hypotheses from \texttt{ConwayAdj} and an involution is \textsc{Pending}.
The two order-3 statements \texttt{z3\_fpf\_orbit\_matrix\_nonexistence} and \texttt{z3\_fixed3\_orbit\_matrix\_nonexistence} contain \texttt{sorry}.
As a result, the summary theorem \texttt{conway\_automorphism\_\allowbreak group\_\allowbreak restricted} is \textsc{Compiled}.

\section{The order-7 proof in Lean}
\label{sec:z7}

\subsection{From an automorphism to an orbit matrix}
\label{sec:z7_reduction}

Let $s$ be a permutation of $V(G)$ with orbits $O_0, \dots, O_{m-1}$ of sizes $n_i = |O_i|$.
For $u \in O_i$, let $B_{ij}$ be the number of neighbors of $u$ in $O_j$.
If $s$ preserves adjacency, then $B_{ij}$ does not depend on the choice of $u$.
In the Lean structure \texttt{Z7OrbitMatrix}, $m = 15$, $n_0 = 1$, and $n_i = 7$ for $i \ge 1$.
The structure also has these conditions for all indices $i$ and $j$:
\begin{enumerate}
 \item $\sum_{j} B_{ij} = 14$, $n_i B_{ij} = n_j B_{ji}$, and $(B^2)_{ij} + B_{ij} = 12\,\delta_{ij} + 2\,n_j$.
 \item $B_{ij} \le n_j$ and $B_{ii} \in \{0, 2\}$. If $i, j \ge 3$, then $B_{ij} \le 4$.
 \item $B_{0j} = 7$ if $j \in \{1, 2\}$, and $B_{0j} = 0$ for the other $j$. $B_{i0} = 1$ if $i \in \{1, 2\}$, and $B_{i0} = 0$ for the other $i$.
 \item $B_{11} = B_{22} = 0$ and $B_{12} = B_{21} = 1$.
 \item The pair $(B_{i1}, B_{i2})$ is $(2, 0)$ for $3 \le i \le 5$, $(0, 2)$ for $6 \le i \le 8$, and $(1, 1)$ for $9 \le i \le 14$.
\end{enumerate}

\begin{theorem}[\texttt{conway\_z7\_induces\_orbit\_matrix}; \textsc{Proved}]
\label{thm:z7_reduction}
Let $A$ satisfy \texttt{ConwayAdj}. Let $s$ be a bijection of the vertex set $\{0, \dots, 98\}$ with $s^7 = \mathrm{id}$, $s \ne \mathrm{id}$, and $A_{s(u)\,s(v)} = A_{uv}$ for all vertices $u$ and $v$. Then a matrix that satisfies \texttt{Z7OrbitMatrix} exists.
\end{theorem}

The Lean statement has no other hypotheses. It does not take an orbit partition, an orbit count, or a quotient matrix as input. The proof has these steps:
\begin{enumerate}
 \item \texttt{Z7Cycles.lean}: each orbit of $s$ has size 1 or 7. The number of fixed vertices is congruent to 99 modulo 7. As a result, a fixed vertex $x_0$ exists.
 \item \texttt{GraphLocalCounts.lean}: \texttt{ConwayAdj} gives degree 14, one common neighbor for adjacent vertices, and two common neighbors for distinct vertices that are not adjacent. The unique common neighbor gives a perfect matching on each neighborhood.
 \item \texttt{Z7FixedGeometry.lean} and \texttt{Z7GraphFrame.lean}: $s$ moves each neighbor of $x_0$. The set $\Gamma(x_0)$ is the union of two orbits of size 7, and no edge joins two vertices of one of these orbits. The vertex $x_0$ is the only fixed vertex.
 \item \texttt{Z7OrbitEnumeration.lean}: a computable list contains one representative of each orbit. The list has 15 entries. Indices 0, 1, and 2 hold $x_0$ and the two orbits in $\Gamma(x_0)$.
 \item \texttt{OrbitQuotient.lean} and \texttt{OrbitReindex.lean}: the matrix $B$ is well defined. It satisfies condition~1 and the bound $B_{ij} \le n_j$. If $n_i$ is odd, then $B_{ii}$ is even.
 \item \texttt{Z7QuotientBoundary.lean} and \texttt{Z7RootedConstruction.lean}: conditions 3 and 4 follow from the structure of $\Gamma(x_0)$.
 \item \texttt{Z7CanonicalArithmetic.lean}: for $i \ge 3$, the quotient equation gives $B_{i1} + B_{i2} = 2$. The equations then give three outer orbits of type $(2, 0)$, three of type $(0, 2)$, and six of type $(1, 1)$. An explicit permutation puts the outer orbits in the order of condition~5.
\end{enumerate}

\subsection{The exterior equations}

The 12 outer orbits have indices $3, \dots, 14$. We write them as $0, \dots, 11$. Let $\ell_i$ be $2$, $0$, or $1$ for the blocks $\{0,1,2\}$, $\{3,4,5\}$, and $\{6, \dots, 11\}$. Let
\begin{equation}
 T_{ij} = 12\,\delta_{ij} + 14 - \ell_i \ell_j - (2 - \ell_i)(2 - \ell_j).
 \label{eq:target}
\end{equation}

\begin{definition}[\texttt{ExteriorEquations}]
A matrix $C \in \mathbb{N}^{12 \times 12}$ satisfies the exterior equations if, for all $i$ and $j$:
$C_{ij} = C_{ji}$; $C_{ij} \le 4$; $C_{ii} \in \{0, 2\}$; $\sum_k C_{ik} = 12$; $C_{i0} + C_{i1} + C_{i2} = C_{i3} + C_{i4} + C_{i5}$; and $\sum_k C_{ik} C_{jk} + C_{ij} = T_{ij}$.
\end{definition}

The theorem \texttt{Z7OrbitMatrix.exterior\_row\_problem} (\textsc{Proved}) shows that the outer $12 \times 12$ block of each \texttt{Z7OrbitMatrix} satisfies the exterior equations.
The theorem \texttt{no\_z7\_orbit\_matrix\_of\_no\_exterior\_equations} (\textsc{Proved}) gives the converse direction that the proof needs. If no matrix satisfies the exterior equations, then no \texttt{Z7OrbitMatrix} exists.

\subsection{Canonical first rows}

A permutation $\sigma$ of $\{0, \dots, 11\}$ that fixes 0 and maps each block to itself keeps the exterior equations true for the matrix $C_{\sigma(i)\sigma(j)}$.
Such a permutation can sort the entries of row 0 inside each block.

\begin{theorem}[\texttt{exterior\_has\_canonical\_root}; \textsc{Proved}]
Let $C$ satisfy the exterior equations. Then a permutation $\sigma$ of this type exists such that $C_{\sigma(i)\sigma(j)}$ satisfies the exterior equations. The row 0 of this matrix is in a fixed list of 33 rows.
\end{theorem}

The list is \texttt{z7RootValues}. The theorem \texttt{z7CanonicalRow\_covered} shows that each sorted row that satisfies the row conditions is in the list. These row conditions are the bound, the diagonal, the row sum, the balance, and $\sum_k C_{0k}^2 + C_{00} = 22$.

\subsection{Five rows with arithmetic proofs}

Five of the 33 rows fail for short arithmetic reasons. Each proof derives \texttt{False} from \texttt{ExteriorEquations} and the fixed row 0. No other hypothesis is necessary.
\begin{itemize}
 \item \textbf{Rows 22 and 26.} The theorem \texttt{exterior\_last\_entry\_diag\_bound} shows $C_{00} + C_{0,11} \le 4$ for each solution. The proof combines rows 0 and 11 and bounds their squared entries. Both rows have $C_{00} = 2$ and $C_{0,11} = 3$.
 \item \textbf{Row 3.} This row is $(0,0,2,0,0,2,1,1,1,1,1,3)$. Let $x_k = C_{11,k}$. The row-sum and balance conditions make $x_6 + \dots + x_{11}$ even. Then the dot product of rows 0 and 11 is even. But the quadratic equation gives $12 - 3 = 9$ for this product.
 \item \textbf{Rows 25 and 31.} For these rows, parity conditions in several rows make all entries of one row even. That row sums to 12, so the sum of the squares of its entries is at least 24. The quadratic equation gives a maximum of 22.
\end{itemize}
The modules are \texttt{Z7RootNormBounds.lean}, \texttt{Z7Root22Arithmetic.lean}, \texttt{Z7Root03Arithmetic.lean}, and \texttt{Z7ComplementParity.lean}.

\subsection{Kernel searches for the other 28 rows}

For the other 28 rows, a finite search in the kernel finds no extension of the fixed row 0. The search has three parts.
\begin{enumerate}
 \item \textbf{Row lists.} For an index $i$, the function \texttt{z7UnaryRows}~$i$ lists all rows that satisfy the conditions on one row: the bound, the diagonal, the row sum, the balance, and the diagonal quadratic equation. A generator makes the list. The generator removes partial rows with a Cauchy--Schwarz bound. The theorem \texttt{cauchy\_step} proves that bound. The theorem \texttt{z7UnaryRows\_of\_row} proves that each row of a solution is in the list. The lists have 7{,}576 rows for $i \le 5$ and 10{,}955 rows for $i \ge 6$.
 \item \textbf{Domains.} Let $a$ be the fixed row 0. The domain of index $i$ contains each row $b$ of \texttt{z7UnaryRows}~$i$ with $b_0 = a_i$ and $\sum_k a_k b_k + a_i = T_{0i}$.
 \item \textbf{Search.} The Boolean function \texttt{noExtension} selects rows for the indices $1, \dots, 11$ in sequence. After each selection, it removes the rows that do not agree with the selected row. It returns \texttt{true} if each branch stops with an empty domain. The theorem \texttt{noExtension\_sound} proves that a \texttt{true} result excludes each complete selection of rows that agree in pairs. The theorem \texttt{exterior\_pinned\_false\_of\_search} connects this result to \texttt{ExteriorEquations}.
\end{enumerate}
For each of the 28 rows, \texttt{decide +kernel} evaluates \texttt{noExtension} to \texttt{true}. This tactic computes in the kernel. It does not use compiled code, so it adds no axiom.
A chain of lists removes the rows in sequence. Each step proves that the first row of a solution is in a shorter list. The last list is empty.

\begin{theorem}[\texttt{no\_exterior\_equations}; \textsc{Proved}]
No matrix $C \in \mathbb{N}^{12 \times 12}$ satisfies the exterior equations.
\end{theorem}

\begin{theorem}[\texttt{conway\_no\_z7\_automorphism}; \textsc{Proved}]
\label{thm:no_z7}
Let $A$ satisfy \texttt{ConwayAdj}. Then no bijection $s$ of $\{0, \dots, 98\}$ satisfies $s^7 = \mathrm{id}$, $s(i) \ne i$ for some $i$, and $A_{s(u)\,s(v)} = A_{uv}$ for all $u$ and $v$.
\end{theorem}

The module \texttt{Z7Main.lean} contains both theorems and \texttt{z7\_orbit\_matrix\_nonexistence}. For all three, \texttt{\#print axioms} reports \texttt{[propext, Quot.sound]}. The proof formalizes the result of Behbahani and Lam~\cite{behbahani2011strongly}. It is not a new mathematical result.

\subsection{Memory limits of kernel computation}
\label{sec:resources}

The first form of the proof computed the 11 row lists again for each of the 28 rows. The kernel keeps its reduction caches for one declaration only. As a result, it did not use a list again in the next declaration. One root row took between 6.8 and 13.1 hours. The time was mostly in the domains for the indices $6, \dots, 11$.

The resident set size that \texttt{ps} reported was near 2.4\,GB. The physical footprint of the process was near 21\,GB on a computer with 16\,GB of memory. The footprint includes compressed and swapped memory. In one domain, the Lean process used 18 minutes of processor time in 2.5 hours. The process waited for paged memory for most of that time. We now measure the physical footprint with \texttt{proc\_pid\_rusage} and stop a job at a set limit.

Three changes keep each check small.
\begin{enumerate}
 \item \textbf{Stored lists.} Each list \texttt{z7UnaryRows}~$i$ is a theorem \texttt{z7UnaryRows\_cached\_$i$}. The theorem equates the generator output to a stored list of integer codes. The proof divides the generator at depth 4 into 114 to 326 small theorems. Each part is a separate \texttt{decide +kernel} check. The longest list proof took 390\,s and 715\,MB.
 \item \textbf{Divided searches.} The search for one root row is a set of theorems. Each theorem covers 8 candidates for row 1. A short lemma joins the parts.
 \item \textbf{A divided bound check.} The proof of \texttt{cauchy\_step} uses one check for each of five values. The single check needed more than 5\,GB.
\end{enumerate}
After these changes, one root row took about 2 minutes.
A full sequential build of the 500 modules of the library had no failures. The largest physical footprint was 2.56\,GB. The build took 3 hours and 3 minutes on a laptop with 16\,GB of memory. Other jobs used the processor during that time.

\subsection{Trusted base and reproduction}

The proof trusts the Lean~4.33.1 kernel and the axioms \texttt{propext} and \texttt{Quot.sound}. The final axiom report does not contain \texttt{Classical.choice}, \texttt{sorryAx}, or \texttt{Lean.ofReduceBool}. A reader can check the result with these commands:
\begin{lstlisting}
lake build Conway.Z7Main
lake env lean Conway/Z7Main.lean    # prints the axioms of the three theorems
\end{lstlisting}

\section{SAT encodings and certified refutations}
\label{sec:sat}

The Lean proof of Section~\ref{sec:z7} does not use the results of this section. These results are independent checks at the formula level. The step from the graph to each CNF is \textsc{Compiled}.

\subsection{Orbit variables and symmetry breaking}

Let $g$ be an automorphism of prime order $p$. The 99 vertices split into orbits of size 1 or $p$. The adjacency between two orbits of size $p$ does not change under cyclic shifts. A circulant Boolean vector $c_{i,j} \in \{0,1\}^p$ gives this adjacency. The entry $c_{i,j,k}$ is 1 if vertex 0 of $\mathcal{O}_i$ is adjacent to vertex $k$ of $\mathcal{O}_j$.
The graph conditions become exact cardinality constraints:
\begin{equation}
 \sum_{u \in \Gamma(x)} A_{u,v} =
 \begin{cases}
  14 & \text{if } x = v, \\
  1 & \text{if } x \sim v, \\
  2 & \text{if } x \not\sim v.
 \end{cases}
\end{equation}
Tseitin variables give the CNF. The exact cardinality constraints use the sequential counters of PySAT (\texttt{EncType.seqcounter}).
Lexicographic leader constraints of Crawford et al.~\cite{crawford1996symmetry} remove symmetric copies. For each generator $\pi$ of the group that keeps the orbit partition, the assignment $X$ must satisfy
\begin{equation}
 X \le_{\text{lex}} \pi(X).
 \label{eq:lex_leader}
\end{equation}
For variables $(x_1, \dots, x_n)$ and their images $(y_1, \dots, y_n)$, auxiliary variables satisfy $e_i \leftrightarrow (x_i = y_i)$, and the encoder adds
\begin{equation}
 (\neg x_i \land y_i) \lor (e_i \land (X_{i+1..n} \le_{\text{lex}} Y_{i+1..n})).
\end{equation}

\subsection{A defect in the encoders and its repair}

In earlier encoders, the variable index 1 was also the Boolean constant \texttt{True}.
The functions \texttt{get\_and\_lit} and \texttt{add\_card\_equals} used the integer \texttt{1} for both meanings.
For the constraint $x_1 + x_2 = 1$, the encoder wrote the unit clause \texttt{[-2]}.
That clause sets $x_2 = 0$ and does not contain $x_1$.
The formula then accepted $(x_1, x_2) = (0, 0)$ and rejected $(0, 1)$.
\begin{example}[The literal-1 defect]
\label{ex:literal1}
Because variable 1 was also the constant \texttt{True}, an exact-1 constraint on $\{x_1, x_2\}$ became the unit clause $\neg x_2$. For this reason, encoder primitives need truth-table tests.
\end{example}
The repaired encoders keep variable indices and Boolean constants apart. Explicit gate managers allocate the variables, and index 0 is an internal sentinel.

The directory \texttt{tests/} contains a deterministic regression suite of 101 tests:
\begin{enumerate}
 \item \textbf{Encoder primitives (20 tests).} Truth tables for conjunctions, cardinality constraints, circulant orbit variables, and lexicographic leader constraints in the order-7 encoder and in the two order-3 encoders.
 \item \textbf{Input audit (11 tests).} A symbolic check of 2{,}394 cardinality constraints against the coordinate model of the involution case, the branch constraints, and the index maps.
 \item \textbf{Run supervisor (8 tests).} Time limits, process control, and limits on the proof size.
 \item \textbf{Orbit pairs (5 tests).} The orbit-pair enumeration for the involution case.
 \item \textbf{Positive controls (32 tests).} The production controls of Section~\ref{sec:positive_controls} on Paley(9), the construction of the Berlekamp--van Lint--Seidel graph, and the clause-check harness. This group also contains older tests of small encoders.
 \item \textbf{Orbit matrices (25 tests).} The quotient-level order-7 encoder on small exhaustive cases and known quotients (10 tests), and its run supervisor (15 tests).
\end{enumerate}
On 3 October 2026, all 101 tests passed. The tests support \textsc{Compiled} status for the current encoder sources. They do not change CNF files from earlier runs, and they do not prove that such files express the graph conditions.

\subsection{Positive controls}
\label{sec:positive_controls}

An \texttt{s UNSATISFIABLE} result cannot show that an encoding accepts all graphs that it must accept. The literal-1 defect made a formula too strong. A positive control applies the same encoder to a known graph with the same $\lambda$ and $\mu$. If the encoder rejects that graph, the control fails.

Two production encoders accept the vertex degree as an argument. The first is for an involution with one fixed vertex (\texttt{build\_z2\_f1\_cnf.py}). The second is for an automorphism of order~3 without fixed points (\texttt{scripts/build\_z3\_canonical\_cnf.py}). The order-3 encoder also accepts the number of orbits. The values $\lambda = 1$ and $\mu = 2$ do not change. With the default arguments, both outputs have the same SHA-256 digests as the outputs of the earlier sources. The default formulas have $666{,}309$ variables and $1{,}508{,}157$ clauses for the involution, and $852{,}324$ variables and $1{,}860{,}024$ clauses for order~3.

The controls use two known graphs with $\lambda = 1$ and $\mu = 2$:
\begin{itemize}
 \item Paley(9), which is $\mathrm{srg}(9, 4, 1, 2)$. For each action, the solver finds a model of the complete formula. The decoded graph satisfies the strongly regular conditions.
 \item The Berlekamp--van Lint--Seidel graph $\mathrm{srg}(243, 22, 1, 2)$~\cite{berlekamp1973strongly}. A script builds it as a Cayley graph on $\mathbb{F}_3^5$ from~\cite{zaw2020deza} and checks $A^2 + A = 20I + 2J$. The map $x \mapsto -x$ is an involution with one fixed vertex. A translation is an automorphism of order~3 without fixed points. The known graph sets the primary variables. The check finds values for the auxiliary variables that make all clauses true: $28{,}257{,}515$ clauses for the involution and $27{,}424{,}360$ clauses for order~3.
\end{itemize}
Negative controls also pass. If one primary value of the Paley(9) assignment changes, the formula becomes unsatisfiable.
These controls are \textsc{Compiled}. They do not prove that an encoder is correct for all inputs. They do not cover the order-7 encoder or the order-3 encoder with three fixed vertices.

\subsection{The canonical order-7 CNF}
\label{sec:order7}

Let $g$ be an automorphism of order 7. Then the vertex set has one fixed vertex $x_0$ and 14 orbits of size 7. Two orbits are in $\Gamma(x_0)$, and 12 orbits are in $\Gamma_2(x_0)$. The canonical CNF uses 498 primary variables: 36 internal orbit valences and 462 circulant variables between orbits. The encoder adds lexicographic leader constraints for the 11 nonidentity elements of a stabilizer. The repaired encoder writes a formula with $177{,}097$ variables and $422{,}508$ clauses.

CaDiCaL 1.9.5~\cite{biere2019cadical} solved the formula. The log records the header \texttt{p cnf 177097 422508}, the verdict \texttt{s UNSATISFIABLE}, exit code 20, $58{,}054.87$\,s of process time, and $119{,}353{,}667$ conflicts. The proof file has $26{,}746{,}488{,}053$ bytes. These lines are an extract from the log:
\begin{lstlisting}
c DRAT 127774798 added clauses 50.10%
c DRAT 127286510 deleted clauses 49.90%
c DRAT 26746488053 bytes (25507.44 MB)
s UNSATISFIABLE
c total process time since initialization:     58054.87    seconds
c exit 20
\end{lstlisting}
\texttt{drat-trim}~\cite{heule2013verifying} checked this CNF against that proof in backward mode. It reported \texttt{s VERIFIED} after $103{,}482.837$\,s:
\begin{lstlisting}
c 378191 of 422508 clauses in core
c 66748552 of 127774799 lemmas in core using 12821576879 resolution steps
c 0 RAT lemmas in core; 62829297 redundant literals in core lemmas
s VERIFIED
c verification time: 103482.837 seconds
\end{lstlisting}

\begin{table}[ht]
\centering
\footnotesize
\setlength{\tabcolsep}{3.5pt}
\caption{Two order-7 encodings. No equivalence between the CP-SAT model and this CNF is proved.}
\label{tab:order7_benchmark}
\begin{tabular}{@{}lllll@{}}
\toprule
Source & Object & Tool & Time & Result \\
\midrule
Thakkar--Severini~\cite{thakkar2026forced} & CP model & CP-SAT & $>48$\,h (14 cores) & \texttt{UNKNOWN} \\
This work (solve) & \texttt{p cnf 177097 422508} & CaDiCaL 1.9.5 & $58{,}054.87$\,s & \texttt{s UNSATISFIABLE} \\
This work (check) & 24.91\,GiB DRAT & \texttt{drat-trim} & $103{,}482.84$\,s & \texttt{s VERIFIED} \\
\bottomrule
\end{tabular}
\end{table}

The SHA-256 digests of the two files are:
\begin{itemize}
 \item \texttt{conway\_z7\_canonical.cnf}:\\
   {\footnotesize\texttt{fc67a85f300151b628ffc5d808b5cb384e69efac2a87c9891aab2c8115e0af47}}
 \item \texttt{proof\_z7\_canonical.drat}:\\
   {\footnotesize\texttt{f312d2a83d85bd11f1c9c2428cb67566ca12b11c70527cfebf0f72572da14a39}}
\end{itemize}
The Zenodo record \href{https://doi.org/10.5281/zenodo.22995101}{10.5281/zenodo.22995101}~\cite{conway99_zenodo} contains both files and the two logs. The earlier record \href{https://doi.org/10.5281/zenodo.22983776}{10.5281/zenodo.22983776} contains the same CNF with a $192{,}218{,}491$-byte proof from an earlier run. That record is not this certificate.
The refutation of \texttt{conway\_z7\_canonical.cnf} is \textsc{Proved}. The step from the graph to this CNF is \textsc{Compiled}. This refutation does not decide the CP-SAT instance of Thakkar and Severini.

\subsection{The exterior-row CNF}
\label{sec:exterior_row}

The script \texttt{scripts/z7\_row\_problem\_cnf.py} writes a direct encoding of the exterior equations. This file is \texttt{z7\_row\_problem.cnf}. It is not the canonical order-7 CNF.
\begin{itemize}
 \item The file has $26{,}166$ variables and $311{,}703$ clauses. Its SHA-256 digest is {\footnotesize\texttt{b746471acd527e52da4706a354fc6fd0d225bb638c121f8b7ec6f965ec5a4ee4}}.
 \item The encoding identifies $C_{ij}$ with $C_{ji}$. It adds no symmetry constraints.
 \item CaDiCaL 1.9.5 reported \texttt{s UNSATISFIABLE} with exit code 20, $9{,}051.98$\,s of real time, and $32{,}331{,}971$ conflicts.
 \item \texttt{drat-trim} read $4{,}750{,}239{,}705$ bytes of proof and reported \texttt{s VERIFIED} after $8{,}879.590$\,s. The core has $201{,}494$ of $311{,}703$ clauses, $12{,}551{,}572$ of $56{,}577{,}564$ lemmas, and 0 RAT lemmas.
\end{itemize}
This refutation is \textsc{Proved} for this CNF. The file and its proof are not in the Zenodo record. Five tests in \texttt{tests/test\_z7\_row\_problem\_cnf.py} run the encoder on small exhaustive instances (\textsc{Compiled}). No proof shows that each solution of the Lean predicate is a model of this CNF.

\subsection{Other order-7 formulas}

A CNF of the 15-orbit quotient conditions has $397{,}296$ variables and $806{,}262$ clauses (\texttt{scripts/orbit\_matrix\_encoding.py}, model \texttt{z7-general}). A CaDiCaL run with a limit of 120\,s stopped with \texttt{c UNKNOWN}. This run is \textsc{Explored}.
An older quotient formula adds restrictions from the Frobenius group of order 21. These restrictions do not follow from an automorphism of order 7. \texttt{drat-trim} verified a refutation of that formula. We make no claim about graphs from it.

\section{The open symmetric cases}
\label{sec:open}

Two cases stay open in the literature: an involution with exactly one fixed vertex, and an automorphism of order~3 without fixed points. This article gives no exclusion of either case.

\subsection{The involution with one fixed vertex}

Let $t \in \Aut(G)$ be an involution with $\Fix(t) = \{x_0\}$. The neighborhood $\Gamma(x_0)$ contains 7 matched pairs. The set $\Gamma_2(x_0)$ contains 84 vertices in 42 orbits $O_0, \dots, O_{41}$ of size 2.
The support of an orbit $O_j = \{u, t(u)\}$ is the set of matched pairs of $\Gamma(x_0)$ that contain neighbors of $u$.
Let $C \in \{0, 1\}^{7 \times 42}$ have $C_{i,j} = 1$ if orbit $O_j$ connects to pair $i$.
Then $C = [C_1 \mid C_1]$, where $C_1 \in \{0,1\}^{7 \times 21}$ is the vertex-edge incidence matrix of $K_7$, and
\begin{equation}
 C C^T = 10 I_7 + 2 J_7.
 \label{eq:cct}
\end{equation}

\begin{proposition}[Column weights]
Each binary matrix $C \in \{0, 1\}^{7 \times 42}$ with $C C^T = 10 I_7 + 2 J_7$ has column weight 2 in each column.
\end{proposition}
\begin{proof}
Let $k_j$ be the weight of column $j$. Then $\sum_j k_j = \Tr(C C^T) = 70 + 14 = 84$, and $\sum_j k_j^2 = \mathbf{1}^T (C C^T) \mathbf{1} = 70 + 98 = 168$. This gives
\begin{equation}
 \sum_{j=1}^{42} (k_j - 2)^2 = 168 - 4 \cdot 84 + 4 \cdot 42 = 0.
\end{equation}
Each term is not negative, so $k_j = 2$ for all $j$.
\end{proof}

The group $W = (\mathbb{Z}_2)^7 \rtimes S_7$ of order $645{,}120$ acts on the signed coordinates of the rooted graph. Fix the orbit $O_0$ with support $\{0, 1\}$. The stabilizer of $O_0$ in $W$ has order 15{,}360. The global sign change fixes each orbit label, so the permutation group on the 42 labels has order 7{,}680. We found these orders by enumeration in the signed coordinate model. Under this stabilizer, the other 41 orbits form three classes:
\begin{itemize}
 \item \textbf{Class A (same support, opposite phase):} size 1, representative $O_{21}$, support $\{0, 1\}$.
 \item \textbf{Class B (support with one common pair):} size 20, representative $O_1$, support $\{0, 2\}$.
 \item \textbf{Class C (disjoint support):} size 20, representative $O_{11}$, support $\{2, 3\}$.
\end{itemize}
An earlier draft gave support $\{1, 6\}$ to $O_1$. That support belongs to $O_{10}$, which is in class B and not in class C. The input-audit tests check the representatives $A=21$, $B=1$, and $C=11$ against the edge list of $K_7$ in the encoder.

\subsection{A solver campaign without a result}

We ran CaDiCaL on Google Cloud Platform (16 cores) on the involution branches and on the order-3 encoding without fixed points. The campaign ran for 12.5 days and stopped on 23 September 2026. The solvers had more than $2.4 \times 10^9$ conflicts in total. They found no satisfying assignment and no refutation. The solver logs are not in this source, so a reader cannot calculate the total again.

Two limits apply. First, the campaign started before the encoder repair of 21 September 2026. Some formulas can contain the literal-1 defect. Second, project records name $O_{10}$ as the representative of class C for a part of the campaign. As a result, it is possible that the campaign does not cover class C.
The campaign is \textsc{Explored}. Conflict counts and proof sizes do not measure the part of the search space that a run removes.

\section{Conclusion}

The status of the results is as follows. The article does not decide existence.
\begin{enumerate}
 \item The Lean library proves that a Conway 99-graph has no automorphism of order~7 (\texttt{conway\_no\_z7\_automorphism}, \textsc{Proved}, axioms \texttt{[propext, Quot.sound]}). This formalizes the known result of Behbahani and Lam. The proof uses no SAT solver and no external certificate.
 \item Stored row lists and divided searches keep each kernel check below 3\,GB. Measurements of the resident set size did not show the real memory use. Measurements of the physical footprint showed it.
 \item The library proves standard structural lemmas and arithmetic refutations (\textsc{Proved}). Graph-level wrappers with assumed counting hypotheses are \textsc{Compiled}. The two order-3 statements contain \texttt{sorry}.
 \item Two CNF refutations, of the canonical order-7 CNF and of the exterior-row CNF, are \textsc{Proved} as formulas. The step from the graph to each CNF is \textsc{Compiled}.
 \item The encoders have a regression suite and positive controls (\textsc{Compiled}). A 12.5-day campaign on the two open cases is \textsc{Explored}.
 \item The involution with one fixed vertex and the automorphism of order~3 without fixed points are \textsc{Pending}.
\end{enumerate}

\paragraph{The rigidity conjecture.}
\begin{conjecture}[Rigidity]
If a strongly regular graph with parameters $(99, 14, 1, 2)$ exists, then $\Aut(G) = \{1\}$.
\end{conjecture}
A search that adds a nontrivial automorphism cannot find a rigid graph. A refutation of the two open cases proves the conjecture. Existence then stays open.

\backmatter

\bmhead{Acknowledgements}
The author thanks the developers of CaDiCaL, \texttt{drat-trim}, PySAT, and Lean~4 for their open-source tools.

\section*{Declarations}

\begin{itemize}
\item \textbf{Funding.} No funding was received for this work.
\item \textbf{Competing interests.} The author has no competing financial or non-financial interests.
\item \textbf{Ethics approval and consent to participate.} Not applicable.
\item \textbf{Consent for publication.} Not applicable.
\item \textbf{Data availability.} The Zenodo record \url{https://doi.org/10.5281/zenodo.22995101} contains the canonical order-7 CNF, its DRAT proof, the CaDiCaL 1.9.5 log, and the \texttt{drat-trim} log. The exterior-row CNF and its proof are not in that record. The files \texttt{tests/production\_positive\_control\_results.json} and \texttt{tests/z7\_general\_run\_result.json} in the code repository record the positive-control output and the bounded quotient-level run, with SHA-256 digests.
\item \textbf{Materials availability.} Not applicable.
\item \textbf{Code availability.} The Lean~4 library, the SAT encoders, the regression tests, and the scripts are at \url{https://github.com/antoniomachuca/conway-99-graph}, tag \texttt{v1.0.0}. The order-7 proof is at commit \texttt{de9a5349552f6ebc3b989e1235f605d2a87b4ab8}. The Zenodo record \url{https://doi.org/10.5281/zenodo.23307118} contains an archive of this tag.
\item \textbf{Author contribution.} Antonio Machuca is the only author. He wrote the Lean library, the SAT encoders, the tests, and the manuscript, and he is responsible for all content.
\item \textbf{Use of large language models.} The author used large language models for code maintenance, test generation, proof work, and LaTeX edits. The Lean kernel, the regression tests, and the retained solver and checker logs check only the statements and files that the text names. The output of a model is not an independent mathematical review.
\end{itemize}

\bibliography{references}

\end{document}
